\documentclass[10pt]{article}
\usepackage[paperwidth=42cm, paperheight=29.7cm, margin=0pt]{geometry}
\usepackage{fontspec}
\usepackage{microtype}
\usepackage[brazil]{babel}
\usepackage{amsmath}
\usepackage{amssymb}
\usepackage{tikz}
\usetikzlibrary{calc, positioning, arrows.meta, shapes.geometric, decorations.pathreplacing, fit, backgrounds}

% === TIPOGRAFIA (Swiss Poster) ===
\setmainfont{Gill Sans}
\setsansfont{Gill Sans}
\newfontfamily\titlefont{Futura}[LetterSpace=12]
\newfontfamily\headerfont{Futura}[LetterSpace=5]

% === PALETA SWISS ===
\definecolor{deepblue}{RGB}{15, 23, 42}     % Sections escuras
\definecolor{coral}{RGB}{239, 68, 68}       % Acento coral
\definecolor{skyblue}{RGB}{14, 165, 233}    % SpO2 / OxyHb
\definecolor{amber}{RGB}{245, 158, 11}      % IR / acento
\definecolor{slate50}{RGB}{248, 250, 252}   % Fundo claro
\definecolor{slate200}{RGB}{226, 232, 240}  % Bordas
\definecolor{slate700}{RGB}{51, 65, 85}     % Texto secundário

\usepackage{xcolor}
\pagestyle{empty}

% === MACRO: Bloco de cor sólida ===
\newcommand{\colorblock}[3]{% {cor bg}{cor texto}{conteúdo}
    \begin{tikzpicture}
        \node[rectangle, fill=#1, text=#2, inner sep=12pt,
              text width=\linewidth-24pt, align=left,
              font=\normalsize] {#3};
    \end{tikzpicture}%
}

\newcommand{\titem}[1]{%
    \noindent\hangindent=1.5em\hangafter=1%
    {\footnotesize$\blacktriangleright$}\hspace{0.5em}#1\par\vspace{3pt}%
}

\begin{document}
\noindent
\begin{tikzpicture}[remember picture, overlay]
    \fill[deepblue] (current page.north west) rectangle ([yshift=-4.5cm]current page.north east);
\end{tikzpicture}

% === HEADER BAR ===
\noindent
\begin{minipage}[c][\dimexpr4.5cm\relax][c]{\textwidth}
    \hspace{2cm}%
    \begin{minipage}[c]{0.6\textwidth}
        {\titlefont\fontsize{42}{48}\selectfont\color{white}\MakeUppercase{Oximetria de Pulso}}\\[6pt]
        {\headerfont\large\color{slate200}Monitorização Contínua da Saturação Arterial de Oxigênio}
    \end{minipage}\hfill
    \begin{minipage}[c]{0.3\textwidth}
        \begin{flushright}
            {\headerfont\color{amber}\Large ONE PAGER ICU}\\[3pt]
            {\small\color{slate200}\textit{Nick Mark, MD} $\cdot$ Tradução}\\[2pt]
            {\footnotesize\color{slate700}onepagericu.com}
        \end{flushright}
    \end{minipage}%
    \hspace{2cm}
\end{minipage}

\vspace{0.6cm}

% === CORPO PRINCIPAL ===
\noindent\hspace{2cm}%
\begin{minipage}[t]{\dimexpr\textwidth-4cm\relax}

% --- PRINCÍPIO (barra colorida) ---
\colorblock{deepblue}{white}{%
    {\headerfont\large\MakeUppercase{Princípio Fisiológico}}\\[6pt]
    A oximetria de pulso explora a absorção diferencial de luz entre a hemoglobina oxigenada (OxyHb) e desoxigenada (DeoxyHb) nos espectros \textbf{vermelho} (660\,nm) e \textbf{infravermelho} (940\,nm). O componente pulsátil (AC) isola a contribuição arterial.
}

\vspace{0.5cm}

% --- GRID DE 2 COLUNAS: DIAGRAMAS ---
\noindent
\begin{minipage}[t]{0.48\textwidth}
    \colorblock{slate50}{deepblue}{%
        {\headerfont\normalsize\color{deepblue}\MakeUppercase{Mecânica do Sensor}}\\[8pt]
        \begin{center}
        \begin{tikzpicture}[scale=1.2, font=\sffamily\small]
            % Dedo técnico
            \draw[thick, slate700, fill=slate50!50] (-2.2, -0.7) .. controls (1, -0.7) and (1.5, -0.5) .. (1.8, 0) .. controls (1.5, 0.5) and (1, 0.7) .. (-2.2, 0.7);
            \draw[slate200] (0.8, 0.5) .. controls (1.5, 0.5) and (1.8, 0.3) .. (1.8, 0) .. controls (1.8, -0.3) and (1.5, -0.5) .. (0.8, -0.5);
            % Sensor
            \fill[deepblue] (-0.8, 1.0) rectangle (1.0, 1.4);
            \fill[deepblue] (-0.8, -1.0) rectangle (1.0, -1.4);
            \node[white, font=\tiny\headerfont] at (0.1, 1.2) {LED};
            \node[white, font=\tiny\headerfont] at (0.1, -1.2) {SENSOR};
            % Raios
            \draw[dashed, coral, -Latex, thick] (0.0, 0.95) -- (0.0, -0.95);
            \draw[dashed, amber, -Latex, thick] (0.4, 0.95) -- (0.4, -0.95);
            \node[right, coral, font=\tiny] at (0.0, 0.5) {660nm};
            \node[right, amber, font=\tiny] at (0.4, -0.5) {940nm};
        \end{tikzpicture}
        \end{center}
    }
\end{minipage}\hfill
\begin{minipage}[t]{0.48\textwidth}
    \colorblock{slate50}{deepblue}{%
        {\headerfont\normalsize\color{deepblue}\MakeUppercase{Espectro de Absorção}}\\[8pt]
        \begin{center}
        \begin{tikzpicture}[scale=0.85, font=\sffamily\small]
            \draw[->, slate700] (0,0) -- (6.2,0) node[right, scale=0.7] {nm};
            \draw[->, slate700] (0,0) -- (0,4.2) node[above, scale=0.7] {Abs};
            \draw[ultra thick, skyblue] (0, 3.5) .. controls (2, 2.8) and (3, 2) .. (5.5, 1) node[right, scale=0.7] {\textbf{OxyHb}};
            \draw[ultra thick, coral] (0, 1) .. controls (1, 2) and (2, 3) .. (3, 3) .. controls (4, 3) and (5, 2.2) .. (5.5, 2.2) node[right, scale=0.7] {\textbf{DeoxyHb}};
            \draw[dotted, deepblue, thick] (0.8, 0) -- (0.8, 4) node[above, scale=0.6] {\textbf{RED}};
            \draw[dotted, deepblue, thick] (4.8, 0) -- (4.8, 4) node[above, scale=0.6] {\textbf{IR}};
        \end{tikzpicture}
        \end{center}
    }
\end{minipage}

\vspace{0.5cm}

% --- GRID 4 COLUNAS: FATORES DE PRECISÃO ---
\colorblock{deepblue}{amber}{{\headerfont\large\MakeUppercase{Fatores que Afetam a Precisão}}}

\vspace{0.3cm}

\noindent
\begin{minipage}[t]{0.24\textwidth}
    \colorblock{coral}{white}{%
        {\headerfont\small\MakeUppercase{Pigmentação}}\\[4pt]
        {\footnotesize
        SpO\textsubscript{2} pode \textbf{superestimar} em pele negra.
        Pacientes negros: \textbf{3$\times$} mais risco de hipoxemia oculta.\\[4pt]
        Esmaltes escuros (azul, verde, preto) interferem no sinal.
        }
    }
\end{minipage}\hfill
\begin{minipage}[t]{0.24\textwidth}
    \colorblock{skyblue}{white}{%
        {\headerfont\small\MakeUppercase{Hemoglobinas}}\\[4pt]
        {\footnotesize
        \textbf{CoHb:} Falsamente normal.\\
        \textbf{MetHb:} Trava em $\sim$85\%.\\
        \textbf{SulfaHb:} Falsa baixa.\\[4pt]
        Co-oximetria (4+ ondas) resolve.
        }
    }
\end{minipage}\hfill
\begin{minipage}[t]{0.24\textwidth}
    \colorblock{amber}{deepblue}{%
        {\headerfont\small\MakeUppercase{Baixo Fluxo}}\\[4pt]
        {\footnotesize
        Choque e vasoconstrição reduzem o sinal pulsátil.\\[4pt]
        ECMO/LVAD: fluxo não-pulsátil torna oximetria inútil.
        }
    }
\end{minipage}\hfill
\begin{minipage}[t]{0.24\textwidth}
    \colorblock{slate700}{white}{%
        {\headerfont\small\MakeUppercase{Técnicas}}\\[4pt]
        {\footnotesize
        \textbf{Lag:} 5--15s periférico.\\
        \textbf{Calibração:} Imprecisa $<$75\%.\\
        \textbf{Corantes:} Falsas quedas.\\
        \textbf{Hiperoxemia:} Não detectada.
        }
    }
\end{minipage}

\vspace{0.5cm}

% --- ONDAS PLETISMOGRÁFICAS (barra inferior) ---
\noindent
\begin{minipage}[t]{0.35\textwidth}
    \colorblock{deepblue}{white}{%
        {\headerfont\normalsize\MakeUppercase{Onda Pletismográfica}}\\[6pt]
        {\footnotesize A morfologia da onda fornece informações sobre tônus vascular. Componentes sistólicos e diastólicos indicam a compliance arterial.}
    }
\end{minipage}\hfill
\begin{minipage}[t]{0.60\textwidth}
    \centering
    \begin{tikzpicture}[scale=0.55, font=\headerfont\scriptsize]
        % Normal
        \node[left, deepblue, font=\headerfont\footnotesize\bfseries] at (-0.5, 0.5) {NORMAL};
        \draw[thick, skyblue] (0,0) .. controls (0.2,0) and (0.4,1) .. (0.5,1) .. controls (0.6,1) and (0.7,0.4) .. (0.8,0.4) .. controls (0.9,0.8) and (1.5,0) .. (2,0)
                               (2,0) .. controls (2.2,0) and (2.4,1) .. (2.5,1) .. controls (2.6,1) and (2.7,0.4) .. (2.8,0.4) .. controls (2.9,0.8) and (3.5,0) .. (4,0);

        % Fraco
        \begin{scope}[shift={(5,0)}]
            \node[left, deepblue, font=\headerfont\footnotesize\bfseries] at (-0.5, 0.15) {FRACO};
            \draw[thick, slate700] (0,0) .. controls (0.2,0) and (0.4,0.3) .. (0.5,0.3) .. controls (0.6,0.3) and (0.7,0.1) .. (0.8,0.1) .. controls (0.9,0.2) and (1.5,0) .. (2,0)
                                    (2,0) .. controls (2.2,0) and (2.4,0.3) .. (2.5,0.3) .. controls (2.6,0.3) and (2.7,0.1) .. (2.8,0.1) .. controls (2.9,0.2) and (3.5,0) .. (4,0);
        \end{scope}

        % Vasoconstrição
        \begin{scope}[shift={(10,0)}]
            \node[left, deepblue, font=\headerfont\footnotesize\bfseries] at (-0.5, 0.5) {VASO$\uparrow$};
            \draw[thick, coral] (0,0) .. controls (0.2,0) and (0.4,1) .. (0.5,1) .. controls (0.6,1) and (0.7,0.7) .. (0.8,0.7) .. controls (0.9,0.9) and (1.5,0) .. (2,0)
                                 (2,0) .. controls (2.2,0) and (2.4,1) .. (2.5,1) .. controls (2.6,1) and (2.7,0.7) .. (2.8,0.7) .. controls (2.9,0.9) and (3.5,0) .. (4,0);
        \end{scope}

        % Vasodilatação
        \begin{scope}[shift={(15,0)}]
            \node[left, deepblue, font=\headerfont\footnotesize\bfseries] at (-0.5, 0.75) {VASO$\downarrow$};
            \draw[thick, amber] (0,0) .. controls (0.2,0) and (0.4,1.5) .. (0.5,1.5) .. controls (0.6,1.5) and (0.7,0.2) .. (0.8,0.2) .. controls (0.9,0.5) and (1.5,0) .. (2,0)
                                 (2,0) .. controls (2.2,0) and (2.4,1.5) .. (2.5,1.5) .. controls (2.6,1.5) and (2.7,0.2) .. (2.8,0.2) .. controls (2.9,0.5) and (3.5,0) .. (4,0);
        \end{scope}
    \end{tikzpicture}
\end{minipage}

\vspace{0.5cm}

% --- RODAPÉ ---
\noindent
\colorblock{deepblue}{slate200}{%
    {\headerfont\small\MakeUppercase{Sumário}} {\footnotesize\color{white}$\cdot$ Sempre correlacione a onda com o pulso. Em dúvida, gasometria arterial (SaO\textsubscript{2}) é o padrão-ouro.} \hfill {\headerfont\footnotesize\color{amber}CRITICAL CARE SERIES $\cdot$ 2025}
}

\end{minipage}% end content minipage

\end{document}
